\section{Introduction}
\label{sec:intro}

Accurate, continuous vehicle volume data is fundamental to modern traffic engineering, urban infrastructure planning, signal timing optimization, and emissions modeling. Traditionally, transportation departments have relied on intrusive roadside sensing technologies, including inductive loop detectors, pneumatic road tubes, and piezoelectric weight-in-motion strips. Although functional, these in-pavement sensors carry prohibitive procurement and installation expenses, disrupt traffic flows during installation, and degrade rapidly under heavy axle loads and surface repaving. Conversely, manual vehicle counting from roadside visual surveys or pre-recorded video logs requires substantial human labor---often demanding upwards of 20 to 30 minutes of focused manual review per hour of video footage---and is prone to observer fatigue and clerical errors.

To address these shortcomings, automated video analysis utilizing computer vision has emerged as an attractive, non-intrusive alternative. Early computer vision systems relied primarily on background subtraction, optical flow fields, or handcrafted edge features to isolate moving vehicles. However, these classical methods deteriorate in real-world operating environments characterized by dynamic illumination shifts, diurnal transitions, long vehicle cast shadows, adverse weather conditions, and low-contrast camera angles.

With the advent of deep convolutional neural networks and real-time object detectors, the ``detect-then-track'' paradigm has become the dominant architectural pattern for camera-based vehicle counting. In a foundational study, Majumder and Wilmot~\cite{majumder2023automated} demonstrated automated directional vehicle counting from pre-recorded traffic video using You Only Look Once (YOLOv3~\cite{redmon2018yolov3}) coupled with a Kalman-filter-based tracking mechanism and a virtual tripwire gate. Their system deployed a user-defined mid-line flanked by two auto-generated parallel gate lines to infer directional flow (entry versus exit) across Baton Rouge commercial corridors, reporting an overall counting accuracy of approximately 90\%. Crucially, their statistical evaluation identified a defining behavioral characteristic of vision-based counting: errors were governed by \textit{systematic undercounting} rather than stochastic overcounting.

Despite these promising initial results, a rigorous analysis of the baseline methodology in~\cite{majumder2023automated} reveals four major theoretical and operational limitations:
\begin{enumerate}
    \item \textbf{Centroid Perspective Parallax}: The baseline system tracks the geometric centroid of each bounding box, $(\frac{x_1+x_2}{2}, \frac{y_1+y_2}{2})$. In typical oblique, elevated closed-circuit television (CCTV) views, the bounding box centroid corresponds to vehicle sheet metal or rooftop elevation rather than the roadway surface. Consequently, tall vehicles (such as transit buses and multi-axle freight trucks) trigger the virtual tripwire prematurely relative to the actual roadway plane.
    \item \textbf{Absence of Modal Classification}: The baseline system aggregates all detected objects into a single generic ``vehicle'' class. Modern highway capacity manuals and pavement design standards require differentiated modal counts---distinguishing passenger cars, motorcycles, commercial trucks, and buses---to assess passenger equivalencies and pavement fatigue.
    \item \textbf{Occlusion-Induced Track Loss and Undercounting}: Standard Simple Online and Realtime Tracking (SORT~\cite{bewley2016sort}) discards bounding box detections falling below fixed confidence thresholds. In dense or multi-lane traffic, vehicles frequently undergo partial visual occlusion, causing the detector to temporarily output low-confidence boxes. As a result, tracks fragment, Kalman filters drift, and vehicles crossing the tripwire while occluded are lost, directly driving the observed undercounting bias.
    \item \textbf{Processing-Time Clock Distortion}: In~\cite{majumder2023automated}, time-interval aggregation (e.g., 5-minute and 15-minute volume bins) was calculated using the host computer's wall-clock execution time rather than the intrinsic video stream timestamps. Because their legacy YOLOv3 deployment required 1.7 hours of processing time per 1 hour of video on an NVIDIA GTX 1060 GPU, tying interval boundaries to wall-clock processing time introduced severe temporal drift whenever inference speed fluctuated.
\end{enumerate}

To resolve these fundamental limitations, this paper presents an advanced, modular video vehicle counting framework named \texttt{vcount}. We modernize the baseline virtual tripwire formulation by introducing geometric ground-contact tracking, multi-class trajectory stabilization, occlusion-tolerant track association, and hardware-independent temporal binning.

The primary contributions of this work are summarized as follows:
\begin{itemize}
    \item \textbf{Perspective-Invariant Ground-Contact Tracking and Directional Gating}: We replace naive centroid tracking with bottom-center tire patch projection ($x_{\text{mid}}, y_{\text{max}}$) paired with a 2D cross-product signed-distance virtual gate geometry. We further incorporate a fallback crossing decision for trajectories initialized between the gate and mid-line, suppressing boundary artifacts.
    \item \textbf{Multi-Class Categorization with Trajectory-Lifespan Majority Voting}: We implement granular modal classification covering passenger cars, motorcycles, buses, and commercial trucks (with custom support for specialized paratransit vehicles). Classification flickering across consecutive frames is resolved through confidence-weighted majority voting accumulated over each vehicle's track lifespan.
    \item \textbf{Occlusion-Resilient Tracking and Hardware-Independent Timestamping}: By embedding ByteTrack's~\cite{zhang2022bytetrack} dual-threshold association into the streaming video pipeline, the system maintains trajectory continuity during severe visual overlaps without the computational penalty of global motion compensation. Furthermore, interval aggregation is anchored strictly to video timestamps ($t = \text{frame} / \text{FPS}$), guaranteeing exact, reproducible volume reporting across varying hardware platforms.
\end{itemize}

The remainder of this paper is organized as follows: Section~\ref{sec:related_work} reviews related work in object detection, multi-object tracking, and virtual detection lines. Section~\ref{sec:method} details the mathematical modeling and pipeline architecture. Section~\ref{sec:evaluation} presents empirical benchmarking on inference speed, resolution trade-offs, and tracking throughput. Section~\ref{sec:discussion} analyzes key engineering trade-offs and threats to validity, and Section~\ref{sec:conclusion} concludes the paper with directions for future work.
